\section{Introduction}
% The need for practical verification algorithms for infinite-state dynamical systems remains due to the continuous expansion and integration of cyber-physical systems and the Internet of Things (IoT) within critical infrastructure. This paper introduces (trajectory-based) state safety certificates, (trajectory-based) state persistence certificates and (trajectory-based) transition persistence certificates to automate verification of safety, persistence, Linear Temporal Logic (LTL) and $\omega$-regular specifications of discrete-time dynamical systems.

Discrete-time dynamical systems are important mathematical models used to describe the characteristics of state changes in control systems. While these models can effectively capture system behavior, their infinite state space presents challenges for verification. Current methods for verifying temporal properties of these systems include the state-triplet method~\cite{wongpiromsarn2015automata} and closure certificates~\cite{murali2024closure}. However, these methods introduce new concerns: (i) The state-triplet method suffers from conservatism; (ii) The closure certificates have low synthesis efficiency. Developing a more efficient approach that reduces conservatism for verifying temporal properties in these systems remains a significant challenge.

% Barrier Certificates for State Invariants.
% A barrier certificate \cite{prajna2004safety} is a real-valued function over the state space that is non-negative over the initial states, negative over the unsafe states, and remains non-negative under transitions. Based on this definition and the principle of structural induction, it follows that the zero superlevel set of the barrier certificate over-approximates the set of reachable states. By combining the zero superlevel set with the definition that the certificate must be negative over the unsafe states, safety is ensured by isolating reachable states from unsafe states. The work presented in \cite{wongpiromsarn2015automata} employs an approach, termed the state-triplet approach, which is based on barrier certificates to disprove the negations of linear temporal logic (LTL) specifications represented using $\omega$-automata. The state-triplet method utilizes barrier certificates to provide separation between consecutive transitions of the $\omega$-automata in a manner that rejects all accepting runs. However, this method may require unrolling the $\omega$-automata and necessarily blocking all simple paths that lead to the acceptance condition. Building on the idea of blocking the paths to accepting states, we propose the state safety certificate, which preserves the original requirement of blocking all simple paths to the acceptance condition while avoiding the need for unrolling and checking all simple paths.

Prior research has focused on formal verification of temporal properties---including safety, reachability, and reach-avoid specifications---within dynamical systems \cite{xue2021reach, vzikelic2023learning, ansaripour2023learning, vzikelic2023compositional, neustroev2025neural}. However, specifications for complex systems often require Linear Temporal Logic (LTL) \cite{pnueli1977temporal} to capture temporal behaviors. Verifying LTL properties in control systems introduces several challenges. First, these properties describe behaviors across infinite execution trajectories \cite{baier2008principles}, typically modeled via $\omega$-automata. When employing proof certificates for verification, the state-triplet method requires unrolling all simple paths and blocking each individually. It is non-trivial to determine the number of certificates needed for synthesis and this method is conservative. Alternatively, closure certificates face the issue of enormous search spaces, as certificates are defined on $\mathcal{X} \times Q \times \mathcal{X} \times Q$. Finally, the infinite state space and nonlinear dynamics of control systems pose major obstacles for traditional verification methods.


% The approach has been extended for verification and synthesis for
% more general dynamical systems [1, 2, 15, 16]. These state-triplet
% approach are bound to suffer from conservatism as the verification
% of a general $\omega-$regular property requires refutation of infinitely
% many visits to some state and that in turn requires a well-founded
% argument [6, 26] over transitive closure of transition relation.

% Closure Certificates for Transition Invariants.
% Chatterjee et al. \cite{chatterjee2024sound} propose a family of witnesses for verifying general LTL formulas on polynomial programs. They employ the concept of a ranking function, which decreases in value when states outside the Büchi specification are visited and remains non-negative throughout the computation. This ultimately ensures that states within the Büchi specification are visited infinitely often, thereby enabling program verification and falsification through this approach. We utilize trajectory-based certificates to demonstrate that elements outside the acceptance condition of the NBA will not be visited infinitely often. This, in turn, proves that states within the acceptance condition will be visited infinitely often, allowing us to use trajectory-based certificates to find counterexample trajectories in deterministic discrete-time dynamical systems.

% The closure certificates \cite{murali2024closure} serve as an analogue to transition invariants \cite{podelski2004complete} and can use simpler template functions to prove safety offering an advantage over classic barrier certificates \cite{prajna2004safety}. Furthermore, persistence closure certificates can verify the persistence specifications of systems while mitigating the conservatism of barrier certificates. By applying persistence closure certificates to a product system, closure certificates can verify general linear-time properties. The approach of utilizing certificates to describe transition invariants for verifying dynamical systems is novel. By integrating with an NBA, which represents the negation of the specification, it proves the accepting states are visited finitely often, thereby demonstrating that all system behaviors are rejected, ultimately verifying LTL and $\omega$-regular properties. However, the LTL closure certificate ($\ref{LTLCC}$) is defined over $\mathcal{X} \times Q \times \mathcal{X} \times Q$, where $\mathcal{X}$ represents the state space of the system and $Q$ denotes the set of states of the automaton corresponding to the negation of the LTL specification. The high-dimensional search space leads to longer synthesis time for the certificate compared to barrier certificates. To mitigate this issue, we propose transition persistence certificates. By combining state safety certificates with transition persistence certificates, we can verify both LTL and $\omega$-regular specifications, and the search space is limited to at most $\mathcal{X} \times Q \times Q$.

% The contributions of the paper are listed next.
% \begin{itemize}
%   \item We propose trajectory-based certificates corresponding to state safety certificates, trajectory-based state persistence certificates and trajectory-based transition persistence certificates. We demonstrate that if all trajectory-based certificates can be found, it can be guaranteed that the system satisfies $\omega$-regular specifications. We have proved that the state persistence certificate is more expressive than the persistence closure certificate, as all trajectory-based state persistence certificates can be computed using the persistence closure certificate. Additionally, we discuss how to utilize trajectory-based certificates to identify counterexample trajectories.
%   \item We propose state safety certificates, state persistence certificates and transition persistence certificates for verification. Using them, we can easily compute all corresponding trajectory-based certificates.
%   \item We present an SMT-based method, applied within a given template function, and an SOS programming method to search for (trajectory-based) state safety certificates, (trajectory-based) state persistence certificates, and (trajectory-based) transition persistence certificates.
%   \item We show that state safety certificates and transition persistence certificates can subsume the state-triplet method for verifying continuous-space dynamical systems against LTL specifications. In addition, we demonstrate that state safety certificates and transition persistence certificates have a smaller search space compared to LTL closure certificates.
%   \item We demonstrate how to use state safety certificates and transition persistence certificates to verify dynamical systems against LTL specifications described by $\omega$-automata through our case studies, and also showcase the efficiency of their synthesis.
% \end{itemize}
% The relevant code of this paper has been uploaded to the repository \cite{resp}.

In this work, we consider the formal verification of Linear Temporal Logic (LTL) properties for discrete-time dynamical systems using proof certificates. As an abstraction-free method, proof certificates have gained popularity in verification and static analysis domains \cite{chakarov2013probabilistic,chatterjee2016algorithmic,chatterjee2024sound,chatterjee2024equivalence}. Murali et al. \cite{murali2024closure} proposed closure certificates, proof certificates defined as real-valued functions over state pairs of dynamical systems, for persistence verification, which can be used to validate systems against LTL properties. Wongpiromsarn et al. \cite{wongpiromsarn2015automata} introduced the state-triplet method, which demonstrates that a system satisfies LTL properties by blocking all reachable paths. Inspired by these two methods, we propose transition safety certificates and transition persistence certificates. Their combined use enables verification of LTL specifications expressed via transition-based Nondeterministic Büchi Automata (NBA)\cite{buchi1966symposium}. Based on these certificates, we present an automated synthesis method using Satisfiability Modulo Theories (SMT). We show that combining these certificates reduces the conservatism of state-triplet methods while requiring smaller search spaces than closure certificates. We validate the effectiveness of our synthesis approach through two case studies.
% Specifically, for Linear Temporal Logic (LTL) verification, we adopt NBA with transition acceptance conditions to express LTL formulas for further enhancing synthesis efficiency.


% These certificates leverage the well-foundedness property on transitive closures to ensure that system trajectories comply with the specifications. The certificate function in this work is defined based on state pairs.
% Regarding certificates to verify system properties, such as safety-related certificates\cite{prajna2004safety}, for discrete systems, Murali et al. proposed closure certificates \cite{murali2024closure} for verifying LTL and $\omega$-regular properties. Wongpiromsarn et al. combined automata and barrier certificates, the state-triplet method, to achieve verification of LTL specifications by isolating non-accepting states from accepting states \cite{wongpiromsarn2015automata}. These methods have provided great inspiration for us in designing certificates.
% % For bounded model checking, it is necessary to calculate the \emph{completeness threshold}\cite{biere2006linear} of LTL. The computation of this value is quite intensive, but without reaching this threshold, no formal guarantees can be provided.
% % Murali et al. also proposed Co-Büchi certificates \cite{murali2023co} which borrow ideas from bounded model checking for specification verification.
% There are many applications of ranking functions in verification of programs \cite{podelski2004complete, shen2013generating, chatterjee2024sound, yao2024mostly, zhu2024breaking, sarita2024syndicate}.
% % , Podelski et al. proved the termination of non-nested program loops by synthesizing linear ranking functions \cite{podelski2004complete}. Shen et al. utilized a mixed numerical-symbolic approach to synthesize ranking functions that can prove the termination of loop programs with polynomial guards and polynomial assignments \cite{shen2013generating}. Chatterjee et al. \cite{chatterjee2024sound} leveraged Büchi automata to propose a new family of ranking functions and a synthesis method for general LTL specifications on imperative programs, thereby verifying system termination, reachability, and safety. Liveness properties for distributed protocols can be reduced to a set of safety properties with the help of ranking functions\cite{yao2024mostly}. Zhu et al. develop a termination analysis that is guaranteed to succeed if a loop (expressed as a formula) admits a (lexicographic) polynomial ranking function\cite{zhu2024breaking}. Sarita et al. proposed a complete synthesis algorithm for ranking functions and invariants based on the termination of specific template programs\cite{sarita2024syndicate}.
% However, there are few applications in the verification of dynamical systems. Persistence closure certificates \cite{murali2024closure} use the decreasing property for system verification. We introduce the notion of ranking functions and combine it with automata to construct specification witnesses.
% % NBA接受条件
% Converting LTL to a TGBA typically results in a smaller automaton compared to converting it to an LGBA \cite{couvreur1999fly}. Degeneralization \cite{giannakopoulou2002states} can transform a TGBA with multiple acceptance conditions into a TBA with a single acceptance condition, an NBA based on the transition acceptance condition. Several tools are available to produce different types of automata from LTL, including Owl \cite{owl}, SPOT \cite{spot}, and so on. We adopt an NBA based on transition acceptance conditions in this paper, which allows for more efficient LTL verification.

\textbf{Contributions.} The main contributions of this paper are:
\begin{itemize}
  \item We propose \emph{transition persistence certificates} and \emph{transition safety certificates} for verifying LTL properties of discrete-time systems with continuous state spaces.
  \item We present an SMT-based synthesis method, achieving 141$\times$ to 909$\times$ speedup over closure certificates while reducing conservatism compared to state-triplet methods.
\end{itemize}

The organization of this paper is as follows: Section 2 covers preliminary knowledge. Section 3 introduces our proof certificates for LTL verification and compares with other methods. Section 4 presents the synthesis method. Section 5 demonstrates the effectiveness through case studies. Section 6 concludes the paper.


\textbf{Related works: } Verification of temporal logic properties via proof certificates has seen significant advancements. To validate discrete-time systems against LTL specifications, Wongpiromsarn et al. \cite{wongpiromsarn2015automata} developed a state-triplet method that uses barrier certificates to block all simple paths from initial states to accepting states in the automaton, requiring enumeration of state triplets $(q_m, q_m', q_m'')$ along these paths. Murali et al. \cite{murali2024closure} introduced closure certificates for persistence verification, which establish disjunctively well-founded transition invariants \cite{podelski2004transition} for dynamical systems. Additionally, substantial research has addressed the verification of specific temporal behaviors such as reachability, safety, and avoidance\cite{chakarov2013probabilistic,chatterjee2016algorithmic,chatterjee2024sound,chatterjee2024equivalence}. Converting LTL specifications into Transition-based Generalized Büchi Automata (TGBA) typically yields smaller automata compared to Labeled Generalized Büchi Automata (LGBA) \cite{couvreur1999fly}. The degeneralization process \cite{giannakopoulou2002states} can convert a TGBA with multiple acceptance conditions into a Transition-based Büchi Automaton (TBA) featuring a single acceptance condition. Various tools exist for generating different automaton types from LTL formulas, such as Owl \cite{owl} and SPOT \cite{spot}.
